\section{Two raw realizations and a quantitative nonlinear consequence}\label{sec:models}
\subsection{Nonlinear hidden-state observations}
\begin{theorem}[Arbitrary hidden dynamics with dominated observations]\label{thm:hidden}
Let $X_t$ be standard Borel, with arbitrary prepared causal hidden-state transitions, and suppose each conditional report law has a density with respect to the same $\sigma_a$, at every hidden history, time and model. Assume the bounded readout loss is continuous in its compact decision variable. Then Theorem~\ref{thm:main} applies, without compactness of $X_t$, a finite-dimensional filter, or contraction. Uniform effective finite-prefix response certificates imply Theorem~\ref{thm:compile}.
\end{theorem}
\begin{proof}
Under a fixed action word, the conditional report law given the observed past is a mixture of the dominated report kernels, and is therefore dominated by $\sigma_{a_t}$. Iterated disintegration on standard Borel spaces gives conditional Radon--Nikodym densities; multiplying them gives a density of the finite report word with respect to $\sigma_{\mathbf a}$. Repeated normalization makes its integral one. This induction also permits hidden transitions depending on past reports; it does not condition on a hidden path as though that path were independent of the reports. The payoff measure is dominated by the report law, so it has an $L^1$ density as well. Dominated convergence gives \eqref{eq:readoutcont}. The raw response hypotheses, rather than any normalized posterior stability, now give the assertions.
\end{proof}
This covers additive nonsingular Gaussian observations of a nonlinear system, with a positive reference density on the full report space. It also covers interval-noise observations whose support moves with the hidden state and parameter. Gaussian tails remain real report events; they are not removed by conditioning. The hidden state can be a function-valued variable in a separable Banach space.

The next result supplies quantitative information beyond effective approximability. Suppose a family on a common bounded metric hidden space has one measurable causal coupling of its raw innovations, valid for every commonly and adaptively selected action, so that until reports first differ
\begin{equation}\label{eq:expanding}
 d_X(X_t^\theta,X_t^{\theta'})\le C_0 L^t s,\qquad L>1,
\end{equation}
when the parameters and preparations are $s$ apart. Suppose the one-call reports can be coupled with mismatch probability at most $C_1(d_X+s)$. On matching reports require the conditional expectation of $\sup_{d\in D}|\ell(d,Z_t)-\ell(d,Z_t')|$ to be at most $C_2(d_X+s)$, conditional on the paired past, states and matched reports. In particular this holds for a hidden target with a uniform pointwise Lipschitz loss. This joint score condition is stronger than closeness of unconditional losses at a fixed decision, which would not control a report-dependent readout. These are consistent raw coupling conditions, not a Lipschitz bound on posterior normalization.

\begin{theorem}[Discount--expansion modulus]\label{thm:holder}
Under \eqref{eq:expanding} and the stated raw report and loss bounds, a constant $C$ independent of $M$ and of the Borel controller satisfies
\begin{equation}\label{eq:holder}
 |J_{\theta,\beta}(\gamma)-J_{\theta',\beta}(\gamma)|
 \le C\begin{cases}
 s,&\beta L<1,\\
 s(1+|\log s|),&\beta L=1,\\
 s^{\log(1/\beta)/\log L},&\beta L>1,
 \end{cases}\qquad 0<s\le\tfrac12.
\end{equation}
The same bounds compare robust values on parameter sets at Hausdorff distance at most $s$ in such a common family. The three orders are sharp in the class of bounded expanding raw experiments. The constants need not be uniform as $\beta\uparrow1$.
\end{theorem}
\begin{proof}
Use the same action and label uniforms until the first report disagreement, and maximal report couplings at each call. While the reports agree, all retained states, actions and readouts agree, even when their rows are discontinuous. The probability of a disagreement by $t$ is at most $C' L^t s$, since the geometric sum of previous raw report errors is of this order. Outside the disagreement event the loss difference has the same bound. Hence the stage risk difference is at most $\min(1,C''L^t s)$.

Sum this bound against $(1-\beta)\beta^{t-1}$. Split at the integer $T$ for which $L^T s$ first becomes of order one. The early sum is bounded by a constant times $s\sum_{t\le T}(\beta L)^{t-1}$; the late sum is at most a constant times $\beta^T$. A geometric sum, or $T$ terms in the critical case, gives \eqref{eq:holder}. Hausdorff comparison follows by comparing each fixed controller on the two sets, then taking infima without changing that controller class.

For sharpness, take $X_0=s\in[0,1]$, $X_{t+1}=\min(1,LX_t)$, reports independent uniform noise, a singleton action and decision, and stage loss $X_t$. The comparison with preparation $0$ is exactly
$(1-\beta)\sum_{t\ge1}\beta^{t-1}\min(1,L^t s)$.
Its early terms and tail have the three matching orders just computed. This is a raw-modulus sharpness witness, not a claim that its memory problem is difficult.
\end{proof}

For example, controlled logistic maps $x\mapsto\theta_a x(1-x)$ on $[0,1]$, with $0\le\theta_a\le4$, satisfy
$|X_t^\theta-X_t^{\theta'}|\le s(4^t-1)/12$ from a common preparation. Reports $Y_t=X_t+U_t$, where $U_t$ is uniform on $[0,w]$ and $w>0$, have total variation at most $|X_t-X'_t|/w$. Squared prediction loss on $[0,1]$ is $2$-Lipschitz in the hidden target. Thus Theorems~\ref{thm:main} and \ref{thm:holder} apply with $L=4$, even in expanding regimes and although report supports move. For a finite prefix the explicit bound
\[
 \Omega_N(s)=(1-\beta)\sum_{t=1}^N\beta^{t-1}
       \min\left\{1,\sum_{r=1}^t\frac{s(4^r-1)}{12w}
                           +\frac{s(4^t-1)}6\right\}
\]
is controller-uniform. Uniform report-bin approximation on a compact enclosing interval has one-call error at most $2h/w$, from the two boundary cells. Coupling it over $t$ calls gives $\min(1,2th/w)$, an alternative to the tensor bound \eqref{eq:EN}. With computable preparation integrals these are effective certificates. No convergence rate for learning $\theta$ is inferred.

\subsection{Noncommuting instruments with changing rank}
\begin{theorem}[Projective and positive quantum observations]\label{thm:quantum}
On a finite-dimensional Hilbert space let a prepared density matrix evolve through trace-preserving quantum instruments with densities
\[
 \mathcal I_{\theta,t}^a(y)(\rho)
       =\sum_{r=1}^R K_{\theta,t,a,r}(y)\rho K_{\theta,t,a,r}(y)^*,
 \quad \int\sum_r K_{\theta,t,a,r}(y)^*K_{\theta,t,a,r}(y)d\sigma_a(y)=I.
\]
No strictly positive lower effect bound is assumed. Let the declared score expectation at a decision be a bounded continuous-in-decision effect evaluated on the post-measurement state. This is one executable terminal audit at the selected scoring call; the unscored trajectory is not subjected to unrecorded earlier audits. Then Theorem~\ref{thm:main} applies to the classical $M$-state observer. Uniform computable instrument approximations and model moduli imply Theorem~\ref{thm:compile}, including projective effects and ranks changing with state, report or parameter.
\end{theorem}
\begin{proof}
For a fixed action word, composition of the unnormalized completely positive maps gives a positive trace-class matrix density on the product report space. Its trace is the report density and integrates to one. Applying the bounded scoring effect gives the payoff density. Continuity in the decision follows before any division by a report probability. Thus zero-probability events and rank changes cause no denominator problem. Each payoff submeasure is defined by its own terminal audit after that report prefix. A joint classical path of outcomes for incompatible quantum audits is neither used nor assumed. If an audit is actually performed before later calls and changes the state, it must instead be included in the intervening physical instrument.

For the effective assertion, replace maps before normalization. If $\|K_r(y)\|,\|K'_r(y)\|\le B$, then, for a positive state $\rho$ of trace one,
\[
 \|K_r\rho K_r^*-K'_r\rho K_r'^*\|_1
 \le 2B\|K_r-K'_r\|.
\]
Summing and integrating gives a raw instrument error bound. Positive trace-preserving instruments, followed by any stochastic controller rows, contract the direct-sum trace norm on Hermitian inputs. Telescoping gives the sum of the one-call errors for every fixed controller. Common parameter and report approximations therefore provide the required moduli, irrespective of posterior rank or contraction. This argument uses neither an invented hidden classical quantum trajectory nor a positive inverse effect.
\end{proof}
A concrete natural instance has two noncommuting measurement actions, with projectors $P_\theta^a$ and $I-P_\theta^a$, and report maps $\rho\mapsto P_\theta^a\rho P_\theta^a$ and its complement. The common reference assigns mass $1/2$ to each of the two reports (so the instrument densities are twice the displayed maps) and dominates all reports, but an outcome can be impossible for some states and parameters. For rotating projectors,
$\|P_\theta^a-P_{\theta'}^a\|\le C|\theta-\theta'|$, so the sum of the two branch errors is at most $4C|\theta-\theta'|$. Report discretization error is zero. Pure preparations and repeated unitary evolution are permitted. This example fails a uniform $\kappa I\le E_y$ hypothesis, yet satisfies the present theorem. The physical density matrix belongs to the experiment; $M$ is classical controller memory, not unrestricted quantum memory.

\begin{corollary}[Singular reference and flagged strata]\label{cor:strata}
The response theorem remains valid for a common Cantor report reference, a countable report reference, or a finite mixture of mutually singular references with observed stratum tags, provided finite histories have the corresponding product densities. Vanishing stratum weights and rank changes are allowed. Effective conditional-expectation certificates use the actual cylinder or tail masses.
\end{corollary}
\begin{proof}
The tensor and averaging proofs use only a finite product of probability spaces and $L^1$ approximation; neither has a Lebesgue-density step. Tagged finite mixtures are probability references on a disjoint union. Zero mixture weights simply set some response densities to zero. For countable reports, retain the tail as a report cell and use its supplied uniform actual mass to bound coarsening; do not condition it away. The quantum construction above supplies genuine changing-rank branches. These verify the hypotheses directly.
\end{proof}
